\section{Measure work inside file services as well as instructions}
Two complete private prototypes now implement the six interfaces. The block version
gathers each channel into windows of at most sixty-four outputs and reuses the earlier
FIR kernel. Its file routines limit each request to 4096 bytes. The direct version
computes each output from the interleaved input and requests the remaining file extent.
Both produce the same canonical output for the worked and listening inputs.

For the measurements below, the file fixture accepts at most seven read bytes and
eleven write bytes per call on the short example. On the two-second input it accepts
at most 2000 bytes per call, which still forces a short transfer from a 4096-byte request.
The exact returned counts determine progress in both programs. File transfer is a byte
operation; a short transfer may split a sample or a frame without changing the final
complete representation.

\tagpdfsetup{table/header-rows={1}}
\begin{center}
\begin{tabular}{llrrr}
Input & Method & Whole MIPS & File bytes & Destination reads \\
\hline
Four frames & Block & 3374 & 60 & 245 \\
Four frames & Direct & 2551 & 60 & 245 \\
Two seconds & Block & 13441424 & 128044 & 260193 \\
Two seconds & Direct & 16577642 & 128044 & 4162969 \\
\end{tabular}
\end{center}

Whole MIPS counts completed instructions from \texttt{main} to
\texttt{results\_ready}, before result publication, using the normal data-address
layout. Destination reads count actual execution-origin byte accesses to the encoded
file arena, including reads performed inside a syscall. The comparison validates both
programs at normal and wide addresses, but these columns are not host elapsed times.
The two programs also use different filtering algorithms, so the whole-program counts
do not isolate request size alone.

The current write service copies its requested bytes into a host buffer before asking
the environment how many it accepts. It also reads the following guard byte. After a
short write, requesting the entire remainder copies many of those same bytes again.
The two-second output takes sixty-five writes in either fixture run. The direct
version requests 4162860 bytes in total; the bounded version requests 260084. Adding
one guard read per write and forty-four header bytes read during save validation gives
the two observed destination totals exactly.

For a fixed small accepted count, summing successive whole remainders grows
quadratically with file length. A request bound limits the copied bytes on each retry.
Here it reduces destination reads by about a factor of sixteen while preserving every
output byte and all file-cleanup rules. Bounded read requests also limit the temporary
buffer requested by the input service. The value 4096 is this prototype's policy, not
a WAV chunk size, a frame stride or a free-writing correctness requirement.

The sample checks also include a separate full-length stereo case with 110274 clipped
results. A sixteen-bit total would return 44738 even if its output samples were all
correct. The averaging input in the table has no clips; do not transfer that observation
to a different coefficient set. The recorded figures describe actual MARS execution
and saved-file bytes. They do not establish IntelliJ launch latency or an IDE capture.
